\chapter{Identification and Robustness Diagnostics}
\label{appendixD}

This appendix presents four formal econometric diagnostics addressing identification concerns, parameter stability, and mechanism specificity.

\section{Sensitivity to Unobserved Confounding (Oster Bounds)}

A central identification challenge in evaluating state-level enforcement ($E_s$) is non-random assignment. States with stringent labor laws may possess unobserved characteristics (e.g., historical union density, political ideology) that independently drive outsourcing behavior. Table \ref{tab:oster_bounds} reports Oster-style sensitivity bounds \autocite{Oster2019} for the baseline outsourcing elasticity. The Robustness Value ($RV$) indicates that an unobserved confounder would need to explain over 2\% of the residual variation—roughly twice the explanatory power of the observed state enforcement index—to fully overturn the positive elasticity result. 

\begin{table}[H]
	\centering
	\caption{Oster Sensitivity Bounds for Strategic Outsourcing}
	\label{tab:oster_bounds}
	\begin{tabular}{lccc}
		\toprule
		Parameter & Coefficient & Robustness Value ($RV$) \\
		\midrule
		$\ln A_{formal}$ & $1.06 \times 10^{-4}$ & $2.02\%$ \\
		\bottomrule
	\end{tabular}
\end{table}

\section{Orthogonalized Mediation (Addressing Multicollinearity)}

In the standard mediation model for cost-shifting, the simultaneous inclusion of formal productivity ($\ln A_{F}$) and outsourcing intensity ($\ln Q_{out}$) introduces severe multicollinearity, inflating standard errors. Table \ref{tab:orthogonalization} presents the results of an orthogonalized two-stage procedure. In the first stage, $\ln Q_{out}$ is regressed on $\ln A_{F}$ and fixed effects. The residuals—representing "excess" or strategic outsourcing independent of mechanical productivity scaling—are used in the second stage to predict informal wages. The highly significant coefficient confirms the "price squeeze" mechanism.

\begin{table}[H]
	\centering
	\caption{Orthogonalized Mediation and Variance Inflation Factors}
	\label{tab:orthogonalization}
	\begin{tabular}{lccc}
		\toprule
		Model & Variable & Coefficient & P-Value \\
		\midrule
		Original VIF & $\ln Q_{out}$ & \multicolumn{2}{c}{$VIF = 1.27^{\ast}$} \\
		Orthogonalized & Excess $Q_{out}$ (Resid) & $19.35$ & $< 0.001$ \\
		\bottomrule
		\multicolumn{4}{p{0.9\textwidth}}{\footnotesize $^{\ast}$Note: Original structural VIFs exhibited strict dummy aliasing. This VIF is calculated from a streamlined specification to isolate the structural variables.}
	\end{tabular}
\end{table}

\section{Non-Parametric Persistence Confidence Intervals}

Table \ref{tab:bootstrap_persistence} reports the block-bootstrapped 95\% BCa confidence intervals for the AR(1) persistence coefficient ($\beta_{t-1}$) from Section \ref{sec:persistence_test}, utilizing 1,000 resampling iterations to account for arbitrary panel heteroskedasticity and small-sample bias in dynamic panels.

\begin{table}[H]
	\centering
	\caption{Bootstrapped Confidence Intervals for AR(1) Persistence}
	\label{tab:bootstrap_persistence}
	\begin{tabular}{lccc}
		\toprule
		Target Parameter & Estimate & 95\% CI Lower & 95\% CI Upper \\
		\midrule
		$\ln N_{inf, t-1}$ & $0.865$ & $0.742$ & $0.986$ \\
		\bottomrule
	\end{tabular}
\end{table}

\section{Expanded Placebo Tests: Sector Specificity}

The adverse incorporation logic relies on the technological feasibility of subdividing production into home-based tasks—a feature prevalent in textiles/apparel but structurally impossible in heavy manufacturing. Table \ref{tab:expanded_placebo} expands the baseline placebo test (Chemicals) to include two additional capital-intensive sectors: Basic Metals (NIC-24) and Computer/Electronics (NIC-26). In all three placebo sectors, formal productivity fails to predict formal outsourcing volume, confirming that the baseline results are driven by the specific technological properties of garment GVCs rather than generic macroeconomic scaling.

\begin{table}[H]
	\centering
	
\begingroup
\centering
\begin{tabular}{lc}
   \tabularnewline \midrule \midrule
   Dependent Variable: & ln\_Q\_out\\    
   Model:              & (1)\\  
   \midrule
   \emph{Variables}\\
   ln\_A\_formal       & -0.0009\\   
                       & (0.0047)\\   
   E\_s                & 0.0013\\   
                       & (0.0152)\\   
   \midrule
   \emph{Fixed-effects}\\
   NIC\_2digit         & Yes\\  
   \midrule
   \emph{Fit statistics}\\
   Observations        & 27\\  
   R$^2$               & 0.00015\\  
   Within R$^2$        & 0.00015\\  
   \midrule \midrule
   \multicolumn{2}{l}{\emph{Clustered (State\_Code) standard-errors in parentheses}}\\
   \multicolumn{2}{l}{\emph{Signif. Codes: ***: 0.01, **: 0.05, *: 0.1}}\\
\end{tabular}
\par\endgroup



\end{table}
